% PLOS-specific author summary; omitted by the PeerJ wrapper.
When biologists group genes by how similarly they are expressed, they usually run
a standard clustering method and then score the result with a separate quality
measure---a measure the method was never trying to do well on. We built a tool,
\tool{}, that removes this gap: the user picks the quality measure they care
about, and a genetic algorithm searches directly for the grouping that scores
best on it. For co-expression data we show that the natural correlation measure
is not an arbitrary choice but the likelihood of a specific picture of what a
module is---a group of genes driven up and down by one shared regulator.

Making that picture explicit lets us ask a question that is usually left
unasked: when is it wrong? We built several artificial datasets that break the
picture in different ways and found that our tool wins on most of them, but fails
badly on one---modules driven by several regulators at once. The failure is
informative rather than fatal, because we can show the search is working
correctly and it is the underlying picture that is wrong: the method finds
groupings that score \emph{better} than the true answer. We then widen the
picture to allow several regulators per module, let the data decide how many are
needed, and recover the lost accuracy. We think reporting both the failure and
its repair is more useful to a reader deciding whether to use the method than a
table in which it always wins.
